\solchapter{particle}

\sol{particle}{1} (a) Forbidden: baryon number ($1\to0$) and lepton number ($0\to-1$) change. Proton decay has never been seen (Problem~\ref{prob:proton-decay}). (b) Allowed: beta decay. (c) Forbidden: muon number ($1\to0$) and electron number ($0\to1$) change separately. (It is allowed in principle by neutrino mixing, but at a rate far too small ever to observe; seeing it would mean new physics.) (d) Allowed: this is how the $\pi^0$ decays, in $10^{-16}\un{s}$. (e) Allowed by the weak interaction, which can change strangeness; it is the $\Lambda$'s main decay, with a lifetime of $2.6\times10^{-10}\un{s}$, long because it is weak. (f) Forbidden by energy and momentum conservation: in the centre-of-mass frame the total momentum is zero, which a single photon cannot have.

\sol{particle}{2} $p = uud$: $\tfrac23+\tfrac23-\tfrac13 = +1$. $n = udd$: $\tfrac23-\tfrac13-\tfrac13 = 0$. $\Delta^{++} = uuu$: $+2$. $\Omega^- = sss$: $-1$. $\pi^+ = u\bar d$: $\tfrac23+\tfrac13 = +1$. $K^- = s\bar u$: $-\tfrac13-\tfrac23 = -1$.

\sol{particle}{3} $\hbar c/m_Wc^2 = 197\un{MeV\,fm}/80\,400\un{MeV} = 2.5\times10^{-3}\un{fm}$, about $1/300$ of a proton's radius. The weak coupling $g\approx0.65$ is actually twice the electromagnetic $e\approx0.3$. But at energies $E\ll m_Wc^2$, the exchanged $W$ is far off its mass shell, and the amplitude is suppressed by about $(E/m_Wc^2)^2$. At LHC energies, the weak force is as strong as electromagnetism, as electroweak unification says it should be.

\sol{particle}{4} $\Gamma(\tau\to e\nu\bar\nu) = \Gamma_\mu(m_\tau/m_\mu)^5$, and this is $17.8\%$ of the total, so
\[
  \tau_\tau = 0.178\,\tau_\mu\left(\frac{m_\mu}{m_\tau}\right)^5 = 0.178\times(2.20\times10^{-6}\un{s})\times\left(\frac{105.7}{1777}\right)^5 = 2.9\times10^{-13}\un{s},
\]
in excellent agreement. In natural units a rate is an energy; $G_F^2$ has dimensions of $\text{energy}^{-4}$, so the rate must be $G_F^2m^5$ times a number. (Physically: the amplitude is $\propto G_F$ and the phase space for three light particles sharing energy $m$ grows as $m^5$.) The agreement checks that the tau couples to the $W$ with the same strength as the muon and electron: \emph{lepton universality}.

\sol{particle}{5} Below charm ($u,d,s$): $N_c(\tfrac49+\tfrac19+\tfrac19) = \tfrac23N_c$. With charm: $\tfrac{10}{9}N_c$. With bottom: $\tfrac{11}{9}N_c$. For $N_c = 3$ these are $2$, $3.33$, and $3.67$, matching the data. (With $N_c = 1$ the predictions would be three times too small.) The measured values are a few percent higher still, because of gluon emission, a correction of $\alpha_s/\pi$.

\sol{particle}{6} The flavour-change probability is maximal when $\Delta m^2c^4L/4\hbar cE = \pi/2$; in practical units $1.27\,\Delta m^2[\mathrm{eV^2}]\,L[\mathrm{km}]/E[\mathrm{GeV}] = \pi/2$, giving $L\approx500\un{km}$. Neutrinos from the sky above have travelled only $10$--$20\un{km}$ from where cosmic rays made them, too short to oscillate. Those from below have crossed the Earth, up to $13\,000\un{km}$, many oscillation lengths; averaged over energies, half of their muon neutrinos have turned into tau neutrinos, which Super-Kamiokande could not see. This up--down asymmetry, measured in 1998, was the discovery of neutrino mass.

\sol{particle}{7} With $m_u\approx2.2\un{MeV}$ and $m_d\approx4.7\un{MeV}$: proton $9\un{MeV}$, neutron $12\un{MeV}$, about $1\%$ of their masses. The other $99\%$ is the energy of the quarks' motion and of the gluon field confining them, as $E = mc^2$ requires. With massless quarks, the proton would still weigh about $900\un{MeV}$. Nearly all the mass of ordinary matter, including you, is QCD binding energy, not Higgs-given mass. (The neutron outweighs the proton partly because $m_d>m_u$, which is why hydrogen is stable.)

\sol{particle}{8} At $91\un{GeV}$: $12\pi/(23\ln(91^2/0.2^2)) = 0.13$ (measured $0.118$; the one-loop formula and the value of $\Lambda$ are approximate). At $2\un{GeV}$: $12\pi/(27\ln100) = 0.30$. As $Q\to\Lambda$, the logarithm goes to zero and $\alpha_s$ blows up: perturbation theory fails and quarks are confined. The positive sign of $33-2n_f$ means the coupling \emph{decreases} at high energy, so quarks behave almost as free particles inside a proton probed at high energy: asymptotic freedom. The $33$ comes from gluons, which carry colour charge themselves and anti-screen; the $-2n_f$ from quark loops, which screen as electrons do in QED. With $17$ flavours, $33 - 34<0$: screening wins, and QCD would behave like QED, weak at low energies and strong at high.

\sol{particle}{9} $y_f = \sqrt2m_f/v$: electron $2.9\times10^{-6}$, bottom $0.024$, top $0.99$. The top quark's coupling is almost exactly $1$, the only fermion mass of the natural size; the others are puzzlingly small. The top is so heavy that it decays (in $5\times10^{-25}\un{s}$) before it can form hadrons, so it is the only quark studied bare, and its loops dominate the quantum corrections to the Higgs mass and the stability of the vacuum.

\sol{particle}{10} The spatial wavefunction is symmetric (all in the ground state), the spin wavefunction $\uparrow\uparrow\uparrow$ is symmetric, and so is the flavour $uuu$: exchanging any two quarks leaves the state unchanged, but identical fermions require it to change sign. Multiplying by the colour wavefunction $\epsilon_{ijk}q_iq_jq_k$, which changes sign when any two indices are swapped, makes the total antisymmetric, as Pauli demands. The combination $\epsilon_{ijk}q_iq_jq_k$ is unchanged by any $SU(3)$ rotation of the colours (it is a determinant, and these rotations have determinant $1$), so it is colour-neutral, ``white''. The only other simple neutral combination is $\bar q^iq_i$, a meson. If only colour-neutral states can exist freely, this explains why hadrons come as $qqq$ and $q\bar q$, and never as free quarks or $qq$ pairs.

\sol{particle}{11} $5\times10^7\un{kg}$ of water is $2.8\times10^9$ moles, or $1.7\times10^{33}$ molecules, each with $10$ protons: $1.7\times10^{34}$ protons. With $\tau = 10^{34}$ years, about $1.7$ decay per year. If no decays are seen in $T$ years, Poisson statistics say the expected number $NT/\tau$ is probably less than about $2.3$ (since $e^{-2.3} = 0.1$, a larger expected number would make seeing none less than $10\%$ likely). So $\tau\gtrsim NT\epsilon/2.3$, with $\epsilon$ the detection efficiency. Super-Kamiokande's limit for $p\to e^+\pi^0$ is over $10^{34}$ years, which has ruled out the simplest grand unified theories.

\sol{particle}{12} \emph{Baryon number violation}: starting from $B = 0$, you cannot reach $B\neq0$ if every process conserves $B$. \emph{C and CP violation}: if C were exact, every process making baryons would be matched by its mirror process making antibaryons at the same rate, and nothing would be gained. The same is true with CP, which also flips the handedness of particles: both must fail, or for every process producing an excess, its conjugate produces an equal deficit. \emph{Departure from equilibrium}: in thermal equilibrium, the population of each state depends only on its energy, and by CPT symmetry a particle and its antiparticle have exactly the same mass. If baryon number is not conserved, there is no chemical potential to favour it either, so equal numbers of baryons and antibaryons result, whatever processes are at work. The expanding universe provides departures from equilibrium (heavy particles decaying out of equilibrium, or phase transitions). The CP violation measured in the Standard Model is about ten orders of magnitude too weak to produce one baryon per billion photons: our existence is evidence of physics we do not yet know.

\sol{particle}{13} At time $t$, the state is $\cos\theta\,e^{-iE_1t/\hbar}\nu_1 + \sin\theta\,e^{-iE_2t/\hbar}\nu_2$, and the amplitude to find $\nu_e$ is $\cos^2\theta\,e^{-iE_1t/\hbar} + \sin^2\theta\,e^{-iE_2t/\hbar}$. Its squared magnitude is
\[
  \cos^4\theta + \sin^4\theta + 2\sin^2\theta\cos^2\theta\cos(\Delta Et/\hbar) = 1 - \sin^22\theta\sin^2(\Delta Et/2\hbar),
\]
using $\cos^4\theta+\sin^4\theta = 1-2\sin^2\theta\cos^2\theta$ and $1-\cos x = 2\sin^2(x/2)$. With $\Delta E = \Delta m^2c^4/2pc\approx\Delta m^2c^4/2E$ and $t = L/c$, the phase is $\Delta m^2c^4L/4\hbar cE$. If $m_1 = m_2$, the two components keep the same relative phase forever and the neutrino stays a $\nu_e$. Oscillation measures differences of squared masses, so seeing it proves that at least one neutrino has mass --- but not how much.
